\documentclass[11pt,fleqn]{report}
\usepackage{graphicx}
\usepackage{amssymb}
\usepackage{float}
\usepackage{longtable}
\usepackage{subcaption}

\usepackage{cml}

\newcommand{\ModelName}{Rocket Motor}
\newcommand{\ModelAuthor}{Daniel Ghan \\ Brian Birmingham}
\newcommand{\ModelAffiliation}{Odyssey Space Research}
\newcommand{\ModelDate}{May 2020}


\begin{document}
\titlePage
\thispagestyle{empty}
\pagenumbering{roman}
\tableofcontents % Print table of contents
\vfill
{\Large \textbf {Revision History}}
{\large
\begin{center}
\begin{tabular}{|l||l|l|l|}
\hline
 \textbf{Version} & \textbf{Date} & \textbf{Author} & \textbf{Purpose} \\ \hline \hline
1 & May 2020 & Daniel Ghan & Initial Version \\ \hline
2 & May 2023 & Brian Birmingham & Requirements and Verification \\ \hline
\end{tabular}
\end{center}
}

\chapter{Introduction} % Make a chapter heading
\pagenumbering{arabic} % Start text with arabic 1
\par The rocket motor models are a set of increasingly complex models of a solid rocket motor. \begin{itemize}
\item \verb|RocketMotor_Basic| is the simplest rocket motor model.  It applies a motor's force and torque, depletes mass from its tank(s), and shuts it off when it depletes its fuel. The thrust and mass flow-rate must be provided directly to the model.
\item \verb|RocketMotor_MassString| allows the user to specify a collection of propellant masses and the type of depletion (parallel consumption or serial consumption).  This option is available in all three models.  
\item \verb|RocketMotor_TableThrust| allows the user to specify the force, specific impulse, mass flow-rate, or total mass consumption as a function of time through the use of look-up tables.
\item \verb|RocketMotor_MultiNozzle| is an extension of \verb|RocketMotor_TableThrust| that supports multiple nozzles. The force and moment generated by the rocket motor is the sum of the forces and moments generated by the nozzles.
\end{itemize}
Note:  The term ``mass'' will always be defined as the mass of the propellant (it could be a single instance of propellant or a group, a \textbf{string}, of propellants).  The mass of the structure or any other component is not considered in this model.

\chapter{Requirements} 
\begin{enumerate}
\item The model shall provide for the computation of forces and moments around a center-of-mass resulting from an idealized thrust vector or vectors.  The center-of-mass is assumed to be that of a rocket body, and the thrust vector(s) are the result of one or more rocket motors.
\item In addition to a structural frame of reference, the model shall provide a motor frame of reference.
\item The default method of thrust characterization is the specification of mass, mass flow-rate, and thrust.  The model shall also be able to accept:
\begin{enumerate}
\item a table of thrust values
\item a specific impule
\item mass burned as a function of time
\end{enumerate}
\item The model shall accept small dispersions in the thrust angles to simulate manufacturing errors.
\item The model shall accept mulit-nozzle configurations, and each nozzle shall be capable of independent operation.
\item The model shall be able to simulate various mass-string configurations and be able to model the propellant flow in a serial or parallel manner.
\end{enumerate}

\chapter{Model Specification}
\section{Code Structure}
\par The rocket motor model consists of 5 classes. The three models summarized in the introduction each have their own classes, with \verb|RocketMotor_Basic| inheriting from CML's \verb|SubscriptionBase|, \verb|RocketMotor_TableThrust| inheriting from \verb|RocketMotor_Basic|, and \verb|RocketMotorMultiNozzle| inheriting from \verb|RocketMotor_TableThrust|. The other two classes are helper classes:   \verb|RocketMotor|-\verb|Dispersion| for dispersing the position and attitude of the motor frame and \verb|RocketMotorNozzle| for modelling the individual nozzles in \verb|RocketMotor_MultiNozzle|.  In addition to the following documentation, the source code has extensive and detailed comments.

\subsection{RocketMotor\_Basic}
\subsubsection{Members}

\begin{center}
\begin{longtable}{|p{2.75cm}|p{2.8cm}|p{7.1cm}|l|} \hline
\textbf{Type} & \textbf{Name} & \textbf{Purpose} & \textbf{Protection} \\
\hline
\endhead
DynamicMass-\newline Group\& & mass\_group & Reference to the dynamic-mass group containing the propellant mass(es) being accessed by the rocket motor. Points either to a dynamic-mass group passed in during construction or to \verb|mass_group_internal|. & protected \\ \hline
DynamicMass-\newline Body* & prop\_mass-\newline \_body & Pointer to the propellant mass-body, used when the propellant is being drawn from a single tank. & protected \\
\hline
DynamicMass-\newline String* & prop\_mass-\newline \_string & Pointer to the propellant mass-string, used when the propellant is being drawn from multiple tanks in a mass-string. & protected \\ \hline
DynamicMass-\newline Body-\newline Properties-\newline Interface\newline & dyn\_mass-\newline \_properties & Reference to the propellant mass properties. & protected \\
\hline
const double\& & time\_now & Reference to simulation time & protected \\
\hline
const double* & veh\_cm & Pointer to position of vehicle centre of mass in structural frame & protected \\
\hline
RocketMotor-\newline Dispersions & dispersions & For dispersing motor frame position and attitude. See section  \ref{subsec:RMD}. & public \\
\hline
bool & commanded & Set to true to command the motor to ignite. & public \\
\hline
bool & force\_mass-\newline \_update & Flag controlling whether to force mass updates to propagate through the mass-tree at the time they are applied versus waiting for a regularly scheduled mass-group update.  If the mass-group is NOT being externally updated (which is likely if the internal mass-group is being used), then this flag should be True or the mass tree will never update. Default: true. & public \\
\hline
bool & motor\_can\_be-\newline \_ shutdown & Determines whether the motor can be shut down by setting \verb|commanded| to false. If true, setting \verb|commanded| to false while the motor is running shuts off the motor. The motor can be re-started as long as there is still mass in the tank. If false, the motor will run until it exhausts its propellant. Default: false. & public \\
\hline
double & thrust-\newline \_magnitude & Magnitude of thrust (input) & public \\
\hline
double[3] & thrust\_unit-\newline \_motor & Thrust direction in motor frame (input) & public \\
\hline
double[3] & position & Position of the rocket nozzle (used to calculate moment) & public \\
\hline
double & mass\_flow\_rate & Rate at which mass is depleted from the tank (input) & public \\
\hline
double[3][3] & T\_struc\_to-\newline \_motor\_ frame & Transformation matrix from the structural frame to the motor frame. Defaults to the Identity matrix.  Modified only in \verb|RocketMotorDispersions::|\newline \verb|apply_dispersions()| & public \\
\hline
double[3] & thrust & The thrust produced by the rocket motor, expressed in the structural frame (output) & public \\
\hline
double[3] & moment & The moment produced by the rocket motor, expressed in the structural frame (output) & public \\
\hline
MotorStatus & status & Current status of the motor. Has three possible values (\verb|MotorStatus| is an enum defined in this class):
\begin{itemize}
\item \textbf{Inactive}: the motor has not yet fired or has been shut down before depleting its propellant
\item \textbf{Firing}: the motor is firing
\item \textbf{Finished}: the motor has exhausted its propellant
\end{itemize}
& protected \\
\hline
DynamicMass-\newline Group & mass\_group-\newline \_internal & An internally-used object for ensuring the changing mass properties are propagated at the same calling rate as this model.  If a mass-group already exists, and the sim is running in multiple threads, this can produce a race condition with the existing group. In this situation, the existing group should be passed in at construction time. & protected \\
\hline
bool & use\_mass-\newline \_string & Indicates whether the motor is drawing mass from a single body or a string of bodies & protected \\
\hline
double & dt & Model cycle time & protected \\
\hline
double & time\_last & Last time update() was executed & protected \\
\hline
double & command\_time & Time at which the ignition command is received & protected \\
\hline
double & burnout\_time & Time of burnout, which tailoff ends & protected \\
\hline
double[3] & thrust\_unit-\newline \_struc & Thrust direction in structural frame & protected \\
\hline
\end{longtable}
\end{center}

\subsubsection{Methods}
\begin{center}
\begin{longtable}{|p{3cm}|p{10cm}|l|} \hline
\textbf{Name} & \textbf{Purpose} & \textbf{Protection} \\
\hline
\endhead
initialize & Initializes the model and disperse the rocket motor's position and attitude & public \\
\hline
update & Updates the model & public \\
\hline
hold\_motor & Temporarily shuts down the motor & public \\
\hline
shutdown\_motor & Permanently shuts down the motor & public \\
\hline
get\_mass & Returns \verb|dyn_mass_properties.consumable_mass| & public \\
\hline
get\_status & returns \verb|status| & public \\
\hline
generate\_torque & Updates \verb|torque| & protected \\
\hline
update\_status & Checks for change in status and executes required actions & protected \\
\hline
update\_mass-\newline\_ consumption & Depletes propellant from the tank & protected \\
\hline
start\_motor & Starts the motor & protected \\
\hline
activate & Activates or re-activates the model & public \\
\hline
\end{longtable}
\end{center}

\subsection{RocketMotor\_TableThrust}
\subsubsection{Members}
\begin{center}
\begin{longtable}{|p{2.75cm}|p{2.8cm}|p{7.1cm}|l|} \hline
\textbf{Type} & \textbf{Name} & \textbf{Purpose} & \textbf{Protection} \\
\hline
\endhead
double & thrust\_max & Maximum thrust produced by the motor. Used to calculate \verb|thrust_fraction|. & protected \\
\hline
double & thrust\_fraction & Current thrust divided by \verb|thrust_max|. For logging purposes. & public \\
\hline
TableLookupSet & table\_set & Interpolation table manager & protected \\
\hline
GenericMulti-\newline InputTable & thrust\_table & Thrust table & protected \\
\hline
GenericMulti-\newline InputTable & isp\_table & Specific-impulse table & protected \\
\hline
GenericMulti-\newline InputTable & mdot\_table & Table for mass burn rate & protected \\
\hline
GenericMulti-\newline InputTable & mburn\_table & Table for total mass burned & protected \\
\hline
Table-\newline Independent-\newline Variable & table\_time & Table of time values (the independent variable) & protected \\
\hline
double & isp & Current specific impulse & protected \\
\hline
double & mburn & Total mass burned so far. If \newline \verb|consumption_type =| \newline \verb|PropConsumptionMburn|, this is obtained from the \verb|mburn_table| and used to calculated the mass flow-rate. & protected \\
\hline
double & delta\_mass & Mass removed on any given frame. Intermediate variable used during calculation of mass flow-rate when \verb|consumption_type =|\newline \verb| PropConsumptionMburn| & protected \\
\hline
double & elapsed\_time & Time since the motor started -- the independent variable for all the look-up tables & protected \\
\hline
Consumption-\newline Type & consumption-\newline \_type & The method used to calculate mass flow-rate. Has five possible values (ConsumptionType is an enum defined in this class): 
\begin{itemize}
\item \textbf{Undefined}: The default value -- will generate an error
\item \textbf{ConstantFlow}: The interpolation tables are not used (the model will act like a \verb|RocketMotor_Basic|)
\item \textbf{PropConsumptionISP}: ISP is obtained from the table and used to calculate mass flow-rate
\item \textbf{PropConsumptionMdot}: Mass \newline flow-rate is obtained from the table
\item \textbf{PropConsumptionMburn}: Total mass burned is obtained from the table and used to calculate mass flow-rate
\end{itemize}
& protected \\
\hline
double & prop\_mass\_init & Initial consumable mass (used to calculate total mass already burned) & protected \\
\hline
static constexpr double & grav\_sea\_level & Standard sea-level gravity, used for ISP calculation & protected \\
\hline
\end{longtable}
\end{center}

\subsubsection{Methods}
\begin{center}
\begin{longtable}{|p{3cm}|p{10cm}|l|} \hline
\textbf{Name} & \textbf{Purpose} & \textbf{Protection} \\
\hline
\endhead
load\_thrust-\newline \_data & Loads data into the thrust table & public \\
\hline
load\_isp\_data & Loads data into the ISP table and performs sanity checks & public \\
\hline
load\_mdot\_data & Loads data into the mass flow-rate table & public \\
\hline
load\_mburn-\newline \_data & Loads data into the mass-burned table & public \\
\hline
load\_time\_data & Loads data into the time table & public \\
\hline
initialize & Initializes the model and the tables & public \\
\hline
update & Updates the model & public \\
\hline
set\_consumption\newline \_type & Switches \verb|consumption_type|. If the model is initialized, checks that the required data have been loaded. & public \\
\hline
shutdown\_motor & Shuts down the motor & public \\
\hline
update\_table & Updates the tables and shuts down the motor if the last time in \verb|table_time| has been reached & protected \\
\hline
compute\_flow-\newline \_ rate\_and\_isp & Compute mass flow-rate and ISP & protected \\
\hline
start\_motor & Starts the motor & protected \\
\hline
\end{longtable}
\end{center}

\subsection{RocketMotor\_MultiNozzle}
\subsubsection{Members}
\begin{center}
\begin{longtable}{|p{2.75cm}|p{2.8cm}|p{7.1cm}|l|} \hline
\textbf{Type} & \textbf{Name} & \textbf{Purpose} & \textbf{Protection} \\
\hline
\endhead
const double\& & atmos\_pressure & Atmospheric pressure, for back-pressure calculations & protected \\
\hline
const double* & motor\_lin\_flex & Linear flex perturbations for all nozzles & protected \\
\hline
const double* & motor\_rot\_flex & Rotational flex perturbations for all nozzles & protected \\
\hline
bool & table\_is\_net-\newline \_ thrust & If true, the thrust-table values represent the magnitude of the sum of the nozzle thrusts; if false, the thrust-table values represent the sum of the magnitudes of the nozzle thrusts. Determines how the model scales the nozzle scale factors during initialization. Already accounts for losses. & public \\
\hline
bool & compute-\newline \_cosine\_ losses & If true, code will compute the public cosine\_loss\_scale\_factor. If false, code will use the input value, which defaults to 1.0. Valid only in the case that \newline \verb|table_is_net_thrust = true|. & public \\
\hline
bool & atm\_press-\newline \_adjust & Adjust the thrust based on the atmospheric pressure & public \\
\hline
double & cosine\_loss-\newline \_scale\_ factor & The magnitude of the sum of the nozzle thrusts divided by the sum of the magnitudes of the nozzle forces. Can be provided by the user or calculated by the model, depending on \verb|compute_cosine_losses|. & public \\
\hline
double & net\_roll\_torq & Copy of \verb|RocketMotor_Basic::moment[0]| & public \\
\hline
double[3] & thrust\_vac & Thrust without back-pressure effect & public \\
\hline
double & thrust\_vac-\newline \_mag & Magnitude of \verb|thrust_vac| & public \\
\hline
bool & using\_flex & Whether flex perturbations are used & protected \\
\hline
double & flex\_threshold & Angle below which rotational flex perturbations are considered negligible. $1.0\times 10^-12$ by default. & protected \\
\hline
size\_t & num\_flex-\newline\_elements & Sanity check passed in at initialization & protected \\
\hline
std::vector\newline  <RocketMotor-\newline Nozzle*> & nozzles\_ptr-\newline \_vec & Vector of pointers to the motor nozzles & protected \\
\hline
size\_t & num\_noz & Number of nozzles on the motor & protected \\
\hline
\end{longtable}
\end{center}

\subsubsection{Methods}
\begin{center}
\begin{longtable}{|p{3cm}|p{10cm}|l|} \hline
\textbf{Name} & \textbf{Purpose} & \textbf{Protection} \\
\hline
\endhead
add\_nozzle & Adds a nozzle to \verb|nozzles_ptr_vec| & public \\
\hline
initialize & Initialize the model & public \\
\hline
force\_initialize & Re-initializes the model & public \\
\hline
initialize\_nozzles & Initializes the nozzles and scales their scale factors based on the number of nozzles and, if \verb|table_is_net_thrust = | \newline \verb|true|, the cosine loss factor & protected \\
\hline
update & Updates the model & public \\
\hline
shutdown\_motor & Shuts down the motor & public \\
\hline
enable\_flex & Enables flex dispersions. If the model is initialized, checks that the required pointers (\verb|motor_lin_flex| and \newline \verb|motor_rot_flex|) have been set. & public \\
\hline
disable\_flex & Disables flex dispersions & public \\
\hline
set\_flex-\newline \_threshold & Sets \verb|flex_threshold| if the new value is valid & public \\
\hline
get\_num\_noz & Returns \verb|num_noz| & public \\
\hline
\end{longtable}
\end{center}

\subsection{RocketMotorNozzle}
\subsubsection{Members}
\begin{center}
\begin{longtable}{|p{2.75cm}|p{2.8cm}|p{7.1cm}|l|} \hline
\textbf{Type} & \textbf{Name} & \textbf{Purpose} & \textbf{Protection} \\
\hline
\endhead
double & azimuth & Nominal nozzle azimuth angle in the motor frame & public \\ 
\hline
double & pitch & Nominal nozzle cant angle in the motor frame & public \\
\hline
double & radius & Nominal distance from nozzle to motor-frame x-axis & public \\
\hline
double & height & Nominal nozzle position along the motor-frame x-axis & public \\
\hline
double & sf & Nominal nozzle scale factor; used to determine how much of the total thrust comes out of this nozzle & public \\
\hline
double & exit\_area & Exit area of nozzle; used for back-pressure calculation & public \\
\hline
double & azimuth\_disp & Azimuth dispersion; affects both position and thrust & public \\
\hline
double & pitch\_disp & Pitch dispersion & public \\
\hline
double & thrust\_azm-\newline \_disp & Thrust azimuth dispersion & public \\
\hline
double & radius\_disp & Radius dispersion & public \\
\hline
double & height\_disp & Height dispersion & public \\
\hline
double & sf\_disp & Scale-factor dispersion & public \\
\hline
double & sf\_true & Dispersed and scaled scale factor & public \\
\hline
double & azimuth\_true & Dispersed azimuth & protected \\
\hline
double & pitch\_true & Dispersed pitch & protected \\
\hline
double & radius\_true & Dispersed radius & protected \\
\hline
double & height\_true & Dispersed height & protected \\
\hline
double[3] & thrust\_dir & Unit thrust direction in the struc-\newline tural frame & public \\
\hline
double[3] & nominal\_thrust-\newline \_dir & Unit thrust direction in the struc-\newline tural frame calculated from nominal values & public \\
\hline
double[3] & position & Nozzle position in structural frame & protected \\
\hline
double & thrust\_mag & Magnitude of nozzle thrust & protected \\
\hline
double & thrust\_vac-\newline \_mag & Magnitude of nozzle thrust without back-pressure effect taken into account & protected \\
\hline
double[3] & thrust & Nozzle thrust & protected \\
\hline
double[3] & thrust\_vac & Nozzle thrust without back-pressure effect taken into account & protected \\
\hline
double & azimuth\_thrust & The azimuth of the thrust direction (may be different from the positional azimuth because of \verb|thrust_azm_disp|) & protected \\
\hline
\end{longtable}
\end{center}

\subsubsection{Methods}
\begin{center}
\begin{tabular}{|p{3cm}|p{10cm}|l|} \hline
\textbf{Name} & \textbf{Purpose} & \textbf{Protection} \\
\hline
shutdown\_nozzle & Shuts down the nozzle & public \\
\hline
initialize & Initialize the nozzle & public \\
\hline
compute\_thrust-\newline \_ mag & Computes the thrust magnitudes & protected \\
\hline
modify\_thrust-\newline \_ mag\_atmos & Applies the back-pressure effect to the thrust magnitude & protected \\
\hline
compute\_thrust-\newline \_ vec & Computes the thrust vectors from the magnitudes and direction & protected \\
\hline
\end{tabular}
\end{center}

\subsection{RocketMotorDispersions}\label{subsec:RMD}
\par \verb|RocketMotorDispersions| contains a single method called \verb|apply_dispersions| that applies dispersions to a position vector and attitude matrix that are passed in as arguments. The members of the class are listed in the table below. (See section \ref{MFD} for further details)
\subsubsection{Members}
\begin{center}
\begin{tabular}{|p{2.75cm}|p{2.8cm}|p{7.1cm}|l|} \hline
\textbf{Type} & \textbf{Name} & \textbf{Purpose} & \textbf{Protection} \\
\hline
double[3] & position-\newline \_dispersion & Dispersion of the position of the motor frame & public \\
\hline
double[3] & motor\_tolerance & Dispersion of the attitude of the motor frame & public \\
\hline
double & tolerance\_mag-\newline \_threshold & Minimum magnitude of \verb|motor_tolerance| below which the attitude will not be dispersed & public \\
\hline
double[3][3] & T\_dispersion-\newline \_to\_ nominal & Transformation from the actual motor\newline frame to the nominal motor frame (output) & public \\
\hline
\end{tabular}
\end{center}

\section{Mathematical Formulation}
\subsection{Reference frames}
\par 	The structural and motor frames are arbitrary, but typically the motor frame is chosen to be the origin of the thrust vector.  Strictly speaking, this model is a ``thrust model'' (not a rocket motor model), so thrusts are specified at a point.  There is, however, a great deal of flexibility in specifying the direction of thrust.  One of the input parameters is the motor unit vector.  This is usually a unit vector with a single non-zero component along the long axis of the structure.  One can choose to ``divide'' the force among any of the three orthogonal unit vector directions, e.g., [0.8, 0.0, 0.6].  Additionally, the motor frame can be offset (dispersed) along the structural frame to simulate mount point deviations.  And finally, manufacturing errors can be introduced by way of motor tolerance angles.  This is discussed further in \ref{MFD}.
\subsection{Thrust, Moment, and Mass Depletion}
\par The single-nozzle models calculate the thrust vector in the motor frame from the thrust magnitude \textit{F} and a unit direction $\hat{u}$. They then transform the vector into the structural frame using the relevant transformation matrix.
$$\mathbf{F} = T_{motor/struc}^T F\mathbf{\hat{u}}$$
\par The moment generated by this thrust about the centre of mass is
$$\mathbf{M} = (\mathbf{x} - \mathbf{x}_{cm}) \times \mathbf{F}$$
where $\mathbf{x}$ and $\mathbf{x}_{cm}$ are the positions of the rocket motor and the vehicle centre of mass, respectively, in the structural frame.
\par The multi-nozzle model calculates the thrust and moment generated by each nozzle in the same way and adds them together.
$$\mathbf{F} = \sum\limits_{i=1}^n \mathbf{F}_i = \sum\limits_{i=1}^n T_{motor/struc}^T F_i \mathbf{\hat{u}_i}$$
$$\mathbf{M} = \sum\limits_{i=1}^n (\mathbf{x}_i - \mathbf{x}_{cm}) \times \mathbf{F}_i$$
The calculation of $F_i$, $\mathbf{\hat{u}_i}$, and $\mathbf{x}_i$ is covered in section \ref{nozzles}.
\par Most often, the model will be supplied a mass flow-rate and calculate the mass lost on each timestep as $\Delta m = \dot{m} \Delta t$. The table-based model has some additional options. The user can supply specific impulse (in seconds), in which case the mass flow-rate is calculated from $$\dot{m} = \frac{F}{I_{sp} g}$$ where \textit{g} is the standard sea-level gravitational acceleration; or the user can supply the total mass burned as a function of time.

\subsection{Motor Frame Dispersion}\label{MFD}
\par During initialization, the model can modify the motor-to-structure transformation matrix to reflect manufacturing errors using the formula below.
$$T'_{motor/struc} = T_{disp}T_{motor/struc}$$
where $T_{disp}$ is a dispersion matrix.  This transformation matrix is generated from Rodrigues' formula, which is based on Euler's theorem (every rotation can be described by an axis of rotation and an angle around it).  Let \verb|RocketMotorDispersions::motor_tolerance| form a vector \textbf{r} of rotation angles.  Rodrigues' formula can then be stated as 
$$T_{disp} = I \cos\theta + \hat{n}_\times \sin\theta + \hat{n}\hat{n}^T (1 - \cos\theta)$$
where 
$$I \equiv  \mbox{Identity matrix} $$
$$\hat{n} =  \frac{\mathbf{r}}{\mathbf{||r||}}\mbox{, the unit vector of rotation}$$
$$\theta = \mathbf{||r||}\mbox{, the angle of rotation in radians}$$
$$\hat{n}_\times =     \left[ \begin{array}{ccc}
0 & -n_z & n_y \\ n_z & 0 & -n_x \\ -n_y & n_x & 0 \end{array} \right] \mbox{, a skew-symmetric matrix from } \hat{n}$$
$$\hat{n}^T \equiv \mbox{the transpose of vector }\hat{n}$$


Although the above equation is not the usual form, it is much faster computationally.  The dispersion calculations are carried out in \verb|RocketMotorDispersions::apply_dispersions()|, and are liberally commented.

\subsection{Multi-Nozzle Motors}
\subsubsection{Nozzle Parameters}\label{nozzles}
\par This section considers a single nozzle. For notational simplicity, the subscript \textit{i} meaning ``of the \textit{i}\textsuperscript{th} nozzle'' is omitted.
\par The nozzle's input parameters are the nominal azimuth $\tilde{\phi}$, pitch $\tilde{\theta}$, radius $\tilde{r}$, height $\tilde{h}$, scale factor $\tilde{\sigma}$, exit area $\tilde{A}$, and the dispersions $\Delta\phi$, $\Delta\phi_F$, $\Delta\theta$, $\Delta r$, $\Delta h$, $\Delta \sigma$, and $\Delta A$. When the nozzle is initialized, the true values of these parameters are computed:
$$\phi = \tilde{\phi} + \Delta\phi$$
$$\theta = \tilde{\theta} + \Delta\theta$$
$$r = \tilde{r} + \Delta r$$
$$h = \tilde{h} + \Delta h$$
$$\sigma = \tilde{\sigma} + \Delta \sigma$$
$\Delta\phi_F$ is the angle between the angle between the horizontal component of the position and the horizontal component of the thrust vector, as shown in \ref{fig:nozzle}.

\begin{figure}[ht]
\includegraphics[scale=0.625]{Images/nozzle_diagram.png}
\caption{Multi-nozzle Geometry}
\label{fig:nozzle}
\end{figure}

\subsubsection{Coordinate System and Assumptions}\label{coord}
\par The nozzles in a multi-nozzle motor are typically arranged in a ring, with the nozzles pointed nearly perpendicular to the plane of the ring. The motor coordinate system is such that the x-axis points away from the primary component of the thrust vector.
\par Each nozzle's position in the motor frame is determined by a height (the distance along the x-axis), a radius (the distance from the x-axis), and an azimuth (measured from the positive z-axis towards the positive y-axis, as shown for Nozzle 1 in the bird's-eye view of \ref{fig:nozzle}):
$$\mathbf{x} = \begin{bmatrix}h \\ r\sin\phi \\ r\cos\phi \\ \end{bmatrix}$$
\par Nominally, the thrust vector is assumed either to intersect with the x-axis or to be parallel to the x-axis. It is determined by the same azimuth from before and pitch, which is defined as the angle between the negative x-axis and the nominal thrust direction, as shown in \ref{fig:nozzle}:
$$\tilde{u} = \begin{bmatrix}-\cos\tilde{\theta} \\ -\sin\tilde{\theta} \sin\tilde{\phi} \\ \sin\tilde{\theta} \cos\tilde{\phi} \\ \end{bmatrix}$$
The true thrust direction is offset from the nominal thrust direction by the pitch dispersion and both azimuth dispersions:
$$\hat{u} = \begin{bmatrix}-\cos\theta \\ -\sin\theta \sin(\phi+\Delta\phi_F) \\ -\sin\theta \cos(\phi+\Delta\phi_F) \\ \end{bmatrix}$$
Note that if $\Delta\phi = \Delta\phi_F = \Delta\theta = 0$, $\tilde{u} = \hat{u}$.

\subsubsection{Division of Force Among Nozzles}
\par In the multi-nozzle model, $F = \verb|thrust_magnitude|$ is not the magnitude of the total thrust vector \textbf{\textit{F}}. Rather, it is used to calculate the nozzle force magnitudes, and the total thrust vector is the sum of the nozzle forces. The force magnitude of each nozzle is
$$F_i = \frac{\sigma_i F}{n}$$
where $n$ is the number of nozzles.
\par When the multi-nozzle model is initialized, it scales the nozzles' scale factors by a ``scaled scale factor'' (\verb|sf_scale|) so that
$$F = \sum\limits_{i=1}^n F_i = \sum\limits_{i=1}^n \frac{\sigma_i F}{n}$$
This equation simplifies to
$$\sum\limits_{i=1}^n \sigma_i = n$$
so \verb|sf_scale| is calculated by dividing the number of nozzles by the sum of the user-set nozzle scale factors (\verb|sf_true|).
\par After initialization, the user is free to change the scale factors to simulate things like failed nozzles. The model does not enforce any bounds on the total force after initialization.
\par Generally, the thrust table will contain raw thrust values--the sum of the magnitudes of the nozzle thrusts--because these values do not depend on the orientations of the nozzles. However, it might contain net thrust values--the magnitude of the total thrust vector. Which convention the model uses is controlled by the \verb|table_is_net_thrust| flag. If the table contains net thrust values, they need to be converted to raw thrust values.
\par Net thrust values would be obtained from raw thrust values be multiplying them by a cosine-loss scale factor:
$$F_{net} = \eta F_{raw}$$
The raw thrust values are recovered from the net thrust values using this equation. The cosine-loss scale factor can be provided by the user or the model will calculate it using the following equation:
$$\eta = \frac{\left|\sum\limits_{i=1}^n\tilde{\sigma}_i \tilde{u}_i \right|}{\sum\limits_{i=1}^n \tilde{\sigma}_i}$$

\subsection{Flex}
\par The multi-nozzle model has an option to ``flex'' the force of each of the nozzles. The user provides pointers to arrays containing the the nozzles' three linear and three rotational flex parameters $\mathbf{f}_{lin}$ and $\mathbf{f}_{rot}$; these can be time-varying values as the model applies them every time-step.
\par The rotational flex parameters are similar to the angles in section \ref{MFD}. They specify an Euler rotation where the angle $\alpha$ is the magnitude of $\mathbf{f}_{rot}$ and the axis $\hat{e}$ is $\mathbf{f}_{rot}$ divided by its magnitude. However, the transformation is applied using a quaternion rather than a matrix. The left transformation quaternion associated with this rotation is (using the scalar-first convention):
$$\mathbf{q} = \begin{bmatrix}
\cos\frac{\alpha}{2} \\[1em]
\sin\frac{\alpha}{2}\frac{\mathbf{f}_{rot}}{\left|f_{rot}\right|} \\
\end{bmatrix}$$
The flexed force $F'_i$ is the unflexed force rotated by this quaternion. The flexed point of application is $\mathbf{x}'_i = \mathbf{x}_i + \mathbf{f}_{lin}$.

\chapter{User's Guide}
\section{Mass setup and considerations}
As described in the Introduction, the rocket motor module comprises three different models.  However, the specification and initialization of the propellant mass is identical across the models.  General guidelines are as follows:
\begin{enumerate}
\item Declare a \verb|MassPropertiesInit| object to represent a mass container (this class is part of the JEOD library).  This object will contain the actual mass value of the propellant, and is the mass of the core body.  It will be treated as a point mass.  This object is used for mass initialization only.  It will not be used in latter stages of the model.  This value should be set in the input file, and will be used as an argument in a later call to \verb|initialize_mass()| (a Trick initialization job).
\item The actual mass variable that will be used for mass, force, and torque calculations is a \verb|DynamicMassBody| object.  This is a CML class, but it inherits from the JEOD class \verb|MassBody|.  The \verb|DynamicMassBody| object invokes an \verb|initialize_mass()| method in a Trick initialization job, but the method is actually a member of \verb|MassBody|.  \verb|initialize_mass()| uses the \verb|MassPropertiesInit| object referenced in item 1.
\item All three motor models can accept the \verb|DynamicMassBody| object in the constructor for that model.  Additionally, all three models can accept a \verb|DynamicMassString| object as the mass representation.  A \verb|DynamicMassString| object is just a collection of \verb|DynamicMassBody| objects with a specified configuration.
\end{enumerate}
Chapter 5 contains several examples of mass setup and initialization.

\section{RocketMotor\_Basic}
\subsection{Instantiation}
\par The \verb|RocketMotor_Basic| constructor has three required arguments and one optional one. The optional argument comes first; it is a manager for the mass depletion. The first required argument can be either a mass body or a mass string; either way, it's what the rocket motor draws its mass from. The second required argument is the simulation time, and the third is a pointer to the center of mass of the motor's parent body.
\begin{code}
class MotorObject : public Trick::SimObject {
  RocketMotor_Basic      basic_motor;
  double                 vehicle_com[3];
  DynamicMassBody        mass_body;
  DynamicMassString      mass_string;
  DynamicMassBodyPropertiesInterface mass_properties;
...
  basic_motor(mass_properties, // optional
              mass_body,       // or mass_string
              jeod_time.time_manager.dyn_time.seconds, 
                              // simulation time
              vehicle_com)    // vehicle center of mass
...
    ("initialization")   basic_motor.initialize();
    (0.005, "effector")  basic_motor.update();
...
};
MotorObject motor_object;
\end{code}

\subsection{Motor Properties}
\par The motor properties are set as follows.
\begin{code}
motor_object.basic_motor.position[0] = 1.0;
motor_object.basic_motor.thrust_direction[0] = 1.0;
// thrust_direction[1] and [2] remain at 0
motor_object.basic_motor.thrust_magnitude = 500000.0;
motor_object.basic_motor.mass_flow_rate = 5.0;
motor_object.basic_motor.dispersions.motor_tolerance[0] =  0.002;
motor_object.basic_motor.dispersions.motor_tolerance[1] = -0.001;
motor_object.basic_motor.dispersions.motor_tolerance[2] =  0.001;
\end{code}
The nozzle position must be specified in the structural frame, but the thrust direction is specified in the motor frame. If the motor frame is not the same as the structural frame, the transformation matrix from the structural frame to the motor frame must be put into \verb|T_struct_to_motor_frame|. Dispersions of the motor frame are automatically applied to the transformation matrix during initialization.  Thrust is constant throughout the entire run.

\section{RocketMotor\_TableThrust}
\subsection{Instantiation}
\par Instantiation of \verb|RocketMotor_TableThrust| is exactly the same as for \verb|RocketMotor_Basic|.  To reduce confusion, this object will be referrred to as a \verb|table_motor|.

\subsection{Motor Properties}
\par The position and orientation of \verb|RocketMotor_TableThrust| are defined in a similar manner to \verb|RocketMotor_Basic|. Thrust and mass flow-rate, on the other hand, have a much more involved set-up. Suppose thrust and mass flow-rate are known as a function of time:

\begin{center}
\begin{tabular}{|c|c|c|}
\hline
\textbf{Time (s)} & \textbf{Thrust (N)} & \textbf{Mass Flow-rate (kg/s)} \\
\hline
0 & 0 & 0 \\
1 & 50 & 2 \\
2 & 150 & 6 \\
3 & 100 & 4 \\
\hline
\end{tabular}
\end{center}

These data can be entered into the model as follows:

\begin{code}
std::vector<double> time_data = {0.0, 1.0, 2.0, 3.0};
std::vector<double> thrust_data = {0.0, 50.0, 150.0, 100.0};
%*\begin{codecomment}
The load\_*\_data functions can accept either a pointer to an array of doubles and the size of the array or a vector of doubles.
\end{codecomment}*)
double flow_rate_data[4] = {0.0, 2.0, 6.0, 4.0};

motor_object.table_motor.load_time_data(time_data);
motor_object.table_motor.load_thrust_data(thrust_data);
motor_object.table_motor.load_mdot_data(flow_rate_data, 4);

%*\begin{codecomment}
Tell the model to use flow-rate data (as opposed to ISP or total mass burned)
\end{codecomment}*)
table_motor.set_consumption_type(
                  RocketMotor_TableThrust::PropConsumptionMdot);
\end{code}

\par The model can also calculate flow-rate if it is given a table of ISP values:
\begin{code}
std::vector<double> isp_data = {20.0, 25.0, 25.0, 15.0};
motor_object.table_motor.load_isp_data(isp_data);
motor_object.table_motor.set_consumption_type(
                  RocketMotor_TableThrust::PropConsumptionIsp);
\end{code}
or a table of the total mass burned as a function of time:
\begin{code}
std::vector<double> total_mass_burned_data = 
                                      {0.0, 1.0, 3.0, 4.5};
motor_object.table_motor.load_mburn_data(total_mass_burned_data);
motor_object.table_motor.set_consumption_type(
                  RocketMotor_TableThrust::PropConsumptionMburn);
\end{code}
or it can be given a constant flow-rate:
\begin{code}
motor_object.table_motor.mass_flow_rate = 3.0;
motor_object.table_motor.set_consumption_type(
                  RocketMotor_TableThrust::Constant);
\end{code}

\section{RocketMotor\_MultiNozzle}
\subsection{Instantiation}
\par \verb|RocketMotor_MultiNozzle| is instantiated similarly to the other two models, but its constructor has another required argument: atmospheric pressure, between simulation time and vehicle center of mass.
\begin{code}
class MotorObject : public Trick::SimObject {
  RocketMotor_MultiNozzle      multi_motor;
  double                       atmospheric_pressure;
  double                       vehicle_com[3];
  DynamicMassBody              mass_body;
  DynamicMassString            mass_string;
  DynamicMassBodyPropertiesInterface mass_properties;
  double                       lin_flex
...
  multi_motor(mass_properties, // optional
              mass_body, // or mass_string
              jeod_time.time_manager.dyn_time.seconds,
                         // simulation time
              atmospheric_pressure, // the name says it all
              vehicle_com)     // vehicle center of mass
...
    ("initialization")   multi_motor.initialize();
    (0.005, "effector")  multi_motor.update();
...
};
MotorObject motor_object;
\end{code}


\subsection{Motor Properties}
\par \verb|RocketMotor_MultiNozzle| inherits all the properties of \verb|RocketMotor_TableThrust|. However, the thrust values in the table may refer to either raw thrust--the sum of the magnitudes of the nozzle thrusts--or net thrust--the magnitude of the nominal total force. By default, the model interprets the table thrusts as raw thrust values; if the table is net thrust, the user must set a flag and either set the cosine-loss scale factor (net thrust divided by raw thrust) or tell the model to calculate it:
\begin{code}
multi_motor.table_is_net_thrust = true;
%*\begin{codecomment} Either \end{codecomment}*)
multi_motor.cosine_loss_scale_factor = 0.9;
%*\begin{codecomment} or \end{codecomment}*)
multi_motor.compute_cosine_losses = true;
\end{code}

\subsection{Nozzles}
\par Nozzles must be added to the motor and their properties must be set before the motor is initialized. Nozzles can be added after initialization, but the motor won't do anything  with them unless it is re-initialized with the \verb|force_initialize()| method. In general, changes made to the nozzles after initialization will have no effect unless the motor is \verb|force_initialize|d.
\begin{code}
RocketMotorNozzle nozzle[2];

for (unsigned int ii = 0; ii < 2; ii++) {
  motor_object.multi_motor.add_nozzle(nozzle[ii]);
  nozzle[ii].exit_area =  3.0; // m2
  nozzle[ii].pitch     =  0.1; // rad
  nozzle[ii].radius    =  0.7; // m
  nozzle[ii].height    = -1.0; // m
}

motor_object.multi_motor.nozzle[0].sf = 1.2;
motor_object.multi_motor.nozzle[1].sf = 0.8;
motor_object.multi_motor.nozzle[0].azimuth = 1.57; // rad
motor_object.multi_motor.nozzle[1].azimuth = 4.71; // rad

// Dispersions
nozzle[0].sf_disp         =  0.001;
nozzle[0].azimuth_disp    =  0.001;
nozzle[0].thrust_azm_disp = -0.003;
nozzle[0].pitch_disp      =  0.002;
nozzle[0].radius_disp     = -0.002;
nozzle[0].height_disp     =  0.003;
nozzle[1].sf_disp         = -0.002;
nozzle[1].azimuth_disp    = -0.004;
nozzle[1].thrust_azm_disp =  0.000;
nozzle[1].pitch_disp      =  0.002;
nozzle[1].radius_disp     =  0.002;
nozzle[1].height_disp     = -0.001;
\end{code}
\par One parameter can be changed after initialization: the nozzle scale factor \verb|sf_true|. This allows the user to simulate things like nozzles failing halfway through the simulation.
\begin{code}
void fail_nozzle_0() {
  // Nozzle 0 fails
  nozzle[0].sf_true = 0.0;
  // This causes an increase in thrust from Nozzle 1
  nozzle[1].sf_true *= 1.5;
}
\end{code}

\subsection{Flex}
\par In order to use flex, the linear and rotational flex parameters must be passed in at initialization:
\begin{code}
  double lin_flex[2][3];
  double rot_flex[2][3];
...  
    ("initialization") multi_motor.initialize(2 * 3,
                                              &lin_flex[0][0],
                                              &rot_flex[0][0]);
\end{code}
The first argument of the initialization method is the length of the flex arrays, which should be three times the number of nozzles. This is purely meant as a sanity check that hopefully prevents the user from accidentally telling the model to read beyond the ends of the arrays.
\par Flex is disabled by default. It can be enabled at any time using the \verb|enable_flex()| method and disabled using \verb|disable_flex|.

\chapter{Verification} 

The model was inspected to ensure compliance with coding standards.  Verification of configurations and algorithms was accomplished through a comprehensive set of unit simulations.  

The verification test suite is grouped into four Trick simulations, with each simulation having distinct run scenarios.  Some runs will produce data that is subsequently checked with known good data.  Other runs will produce errors.  These errors are intentional and are designed to test the model's ability to recognize inappropriate and/or illegal configurations, incorrect user input files, and error conditions that could occur during initialization or run-time.  And a small group of runs are deliberately designed to fail.  These failure conditions will not permit the simulation to run and are reported accordingly.

\section{Basic Rocket Motor Verification}
The run scenarios that test the standard rocket motor configuration, \textbf{SIM\_basic\_rocket\_motor}, are as follows:
\begin{itemize}
  \item RUN\_01\_basic (see \ref{sec:Basic1})
  \item RUN\_02\_not\_active (see \ref{sec:Basic2})
  \item RUN\_03\_restart\_exhaustfuel\_restart (see \ref{sec:Basic3})
  \item RUN\_04\_motor\_thrust (see \ref{sec:Basic4})
  \item RUN\_11\_non\_unit\_direction (see \ref{sec:Basic11})
  \item RUN\_12\_disperse\_NULL (see \ref{sec:Basic12})
  \item RUN\_13\_duplicate\_update (see \ref{sec:Basic13})
  \item RUN\_14\_bad\_thrust\_unit\_vec (see \ref{sec:Basic14})
  \item RUN\_20\_negative\_mass\_FAIL (see \ref{sec:Basic20})
  \item RUN\_21\_no\_cm\_specified\_FAIL (see \ref{sec:Basic21})
\end{itemize}

\setlength\parindent{24pt}

\subsection{Basic}\label{sec:Basic1}
This scenario tests the basic functionality of the single-nozzle rocket motor model.  The required inputs are vehicle center-of-mass, mass of propellant (recall that other masses - tank mass, structure mass, etc. - are not considered), location and direction of nozzle, thrust magnitude, and mass flow rate.  The output data show that the propellant mass depletes per schedule, and that the thrust magnitude and direction is consistent with expectations.  In the first figure below, each data point represents the propellant mass at discrete simulation times.  Since we have constant thrust, a linear reduction is expected and can be seen.  The second figure shows a constant thrust over time, adjusted for the frame-of-reference change from the motor to the structure:

\begin{figure}[h!]
	\centering
	\begin{subfigure}[b]{0.45\linewidth}
		\includegraphics[width=\linewidth]{Images/Basic01.png}
		\caption{Mass decrease}
	\end{subfigure}
	\begin{subfigure}[b]{0.45\linewidth}
		\includegraphics[width=\linewidth]{Images/Basic02.png}
		\caption{Thrust}
	\end{subfigure}
	\caption{Single-nozzle scenario}
\end{figure}
			
\subsection{Not Active}\label{sec:Basic2}
The motor must be commanded to ignite.  This does not have to happen at any particular time, but the \textbf{RocketMotor\_Basic} member variable \textbf{commanded} must be set equal to True in order to start the modelling process.  Otherwise, mass depletion remains at zero, and the thrust and moment vectors are not computed.  This runs sets \textbf{commanded = False} for the entirety of the run.  As expected, no computations occur.

\subsection{Restart, Exhaust Fuel, Restart}\label{sec:Basic3}
As long as the \textbf{RocketMotor\_Basic} member variable \textbf{motor\_can\_be\_shutdown} is set equal to True, the motor can be commanded to ignite and shutdown at will.  In this run, the motor ignites at the beginning of the simulation and fires until the simulation time reaches four seconds.  It is commanded to shut down for two seconds, then restart.  The mass has been chosen so that the fuel is exhausted after eight seconds.  This causes an automatic shutdown.  The motor is then commanded to restart.  No thrust occurs, which is to be expected.  This run exercises various sanity checks in the model, such as no thrust when the motor is off and no thrust when the propellant is exhausted.  The following two plots show the mass and thrust as functions of time.  The motor start/stop cycle can be clearly seen.
\begin{figure}[h!]
	\centering
	\begin{subfigure}[b]{0.45\linewidth}
		\includegraphics[width=\linewidth]{Images/Restart01.png}
		\caption{Mass decrease}
	\end{subfigure}
	\begin{subfigure}[b]{0.45\linewidth}
		\includegraphics[width=\linewidth]{Images/Restart02.png}
		\caption{Thrust}
	\end{subfigure}
	\caption{Stop/Start scenario}
\end{figure}



\subsection{Motor Thrust}\label{sec:Basic4}
The model can accomodate dispersions in a) the position of the motor frame, and b) the attitude of the motor frame.  This unit simulation tests the position dispersion.  Below are two runs...a run with no position dispersion, and a run with a non-zero dispersion in all three dimensions.  The thrust vector will not change, but as expected, the moment vector shows the effects of the dispersion.
\begin{figure}[H]
	\centering
	\includegraphics[width=0.5\linewidth]{Images/Thrust01.png}
	\caption{Moment components}
\end{figure}

\subsection{Non-unit Direction}\label{sec:Basic11}
One must specifiy the thrust direction in the motor frame.  There is no default value.  Additionally, the magnitude of this vector must be unity.    This run checks that condition.
\begin{code}
//  This vector has magnitude not equal 1
rocket_motor.basic_motor.thrust_unit_motor = [1.0, 2.0, 3.0]
}
\end{code}
will result in
\begin{figure}[ht!]
	\centering
	\includegraphics[width=\linewidth]{Images/Nonunit01.png}
	\caption{Non-unit motor thrust direction error}
\end{figure}

\subsection{NULL Dispersion}\label{sec:Basic12}
The position vector must be non-null.  Recall that this is the position of action for the motor expressed in the structural frame.  The transformation matrix from the structural frame to the motor frame is computed from this information plus any dispersions that may exist.  A null position vector will result in a null transformation matrix.  The results of a null position vector and transformation matrix are below.
\begin{figure}[ht!]
	\centering
	\includegraphics[width=\linewidth]{Images/NullDispersion01.png}
	\caption{Null position vector, Null transformation matrix}
\end{figure}

\subsection{Duplicate Update}\label{sec:Basic13}
This run forces a call the to update() method of the model at simulation time zero.
\begin{figure}[ht!]
	\centering
	\includegraphics[width=0.4\linewidth]{Images/DuplicateUpdate01.png}
	\caption{Trick update() at sim time = 0}
\end{figure}


\subsection{Bad Thrust Unit Vector}\label{sec:Basic14}
As mentioned in \ref{sec:Basic11}, the thrust direction must be specified.  A zero thrust direction or an omission will result in a default value of {1.0, 0.0, 0.0}.
\begin{code}
rocket_motor.basic_motor.thrust_unit_motor = [0.0, 0.0, 0.0]
}
\end{code}
will result in
\begin{figure}[ht!]
	\centering
	\includegraphics[width=\linewidth]{Images/BadThrustUnitVector01.png}
	\caption{Illegal thrust unit vector}
\end{figure}

\subsection{Negative Mass (Fail)}\label{sec:Basic20}
Negative propellant mass is not allowed.  This run tests various checks for this condition.  The input.py file for \textbf{RUN\_20\_negative\_mass\_FAIL} has a detailed explanation of the checks.
\begin{code}
rocket_motor.mass_init.mass =0.0
rocket_motor.motor_prop.residual_mass = 2
rocket_motor.motor_prop.dynamic_properties.consumable_mass = -1
\end{code}
results in 
\begin{figure}[ht!]
	\centering
	\includegraphics[width=\linewidth]{Images/NegativeMass01.png}
	\caption{Negative mass}
\end{figure}

\newpage

\subsection{No Center-of-Mass Specified (Fail)}\label{sec:Basic21}
The structural center-of-mass must be specified.  If not, below is the failure message.
\begin{figure}[ht!]
	\centering
	\includegraphics[width=\linewidth]{Images/NoCOM01.png}
	\caption{No center-of-mass specification}
\end{figure}


\section{Mass String Verification}
The run scenarios that test the Mass String capability, \textbf{SIM\_basic\_rocket\_motor\_MassString}, are as follows:
\begin{itemize}
  \item RUN\_01\_parallel (see \ref{sec:String1})
  \item RUN\_02\_series (see \ref{sec:String2})
  \item RUN\_20\_Fail (see \ref{sec:String20})
\end{itemize}

\subsection{Parallel}\label{sec:String1}
When specifying multiple fuel tanks, the propellant can be drawn in a ``serial'' manner (a single tank supplies the propellant until exhausted, then the next tank is selected as the source, and so on) or in a ``parallel'' manner (propellant is drawn from all tanks simultaneously).  This can be configured by setting the \textbf{DynamicMassString} boolean member variable \textbf{flow\_down}.  The parallel configuration sees this variable set to False.  This will deplete mass from all bodies equally.  A mass flow rate of 3 kg/sec was specified.  Note how a depletion rate of 1.5 kg/sec (3 kg/sec divided by the two tanks) is shared equally.
\begin{figure}[H]
	\centering
	\begin{subfigure}[b]{0.45\linewidth}
		\includegraphics[width=\linewidth]{Images/Parallel01.png}
		\caption{Tank A depletion}
	\end{subfigure}
	\begin{subfigure}[b]{0.45\linewidth}
		\includegraphics[width=\linewidth]{Images/Parallel02.png}
		\caption{Tank B depletion}
	\end{subfigure}
	\caption{Parallel tank depletion}
\end{figure}




\subsection{Serial}\label{sec:String2}
Similar to \ref{sec:String1}, the mass depletion will occur in a sequential manner if \textbf{flow\_down=True}, e.g., mass\_body[0] is emptied, then mass\_body[1] is emptied, and so on.  Note how the entire mass flow rate of 3 kg/sec is seen by a single tank.
\begin{figure}[h!]
	\centering
	\begin{subfigure}[b]{0.45\linewidth}
		\includegraphics[width=\linewidth]{Images/Series01.png}
		\caption{Tank A depletion}
	\end{subfigure}
	\begin{subfigure}[b]{0.45\linewidth}
		\includegraphics[width=\linewidth]{Images/Series02.png}
		\caption{Tank B depletion}
	\end{subfigure}
	\caption{Series tank depletion}
\end{figure}

\newpage

\subsection{No mass body}\label{sec:String20}
When using the \textbf{DynamicMassString} class (as in this section), at least one mass body must be added.  Failure to do so results in the following output:
\begin{figure}[ht!]
	\centering
	\includegraphics[width=\linewidth]{Images/NoMassString01.png}
	\caption{No mass body specified}
\end{figure}



\section{Multi-Nozzle Verification}
The run scenarios that test the Multiple Nozzle capability, \textbf{SIM\_multi\_noz\_rocket\_motor}, are as follows:
\begin{itemize}
  \item RUN\_01\_basic (see \ref{sec:MultiNoz1})
  \item RUN\_02\_thrust\_is\_net (see \ref{sec:MultiNoz2})
  \item RUN\_03\_compute\_cosine\_losses (see \ref{sec:MultiNoz3})
  \item RUN\_04\_atm\_pressure (see \ref{sec:MultiNoz4})
  \item RUN\_05\_isp\_table (see \ref{sec:MultiNoz5})
  \item RUN\_06\_mdot\_table (see \ref{sec:MultiNoz6})
  \item RUN\_07\_small\_tolerance (see \ref{sec:MultiNoz7})
  \item RUN\_09\_broken\_nozzle (see \ref{sec:MultiNoz9})
  \item RUN\_10\_scale\_factor\_warnings (see \ref{sec:MultiNoz10})
  \item RUN\_11\_bad\_exit\_area (see \ref{sec:MultiNoz11})
  \item RUN\_12\_consumption\_type\_error (see \ref{sec:MultiNoz12})
  \item RUN\_13\_bad\_update\_call (see \ref{sec:MultiNoz13})
  \item RUN\_14\_duplicate\_nozzle (see \ref{sec:MultiNoz14})
  \item RUN\_15\_negative\_thrust (see \ref{sec:MultiNoz15})
  \item RUN\_20\_init\_warnings\_FAIL (see \ref{sec:MultiNoz20})
  \item RUN\_22a\_NULL\_flex\_FAIL (see \ref{sec:MultiNoz22a})
  \item RUN\_22b\_NULL\_flex\_FAIL (see \ref{sec:MultiNoz22b})
  \item RUN\_23\_flex\_size\_error\_FAIL (see \ref{sec:MultiNoz23})
  \item RUN\_24\_no\_nozzles\_FAIL (see \ref{sec:MultiNoz24})
  \item RUN\_25a\_noz\_init\_FAIL (see \ref{sec:MultiNoz25a})
  \item RUN\_25b\_noz\_init\_FAIL (see \ref{sec:MultiNoz25b})

\end{itemize}

\subsection{Basic}\label{sec:MultiNoz1}
This tests the basic functionality of the multi-nozzle rocket motor model.  It serves as a baseline and building block for the subsequent tests.  The thrust profile can be found in \textbf{example\_rocket\_motor.hh.}  Additionally, it tests the linear and rotational flex modeling capabilities.
Flexing presents itself as a loss of thrust.  The following linear and rotational flex parameters were added to the nozzle configuration:
\begin{code}
rocket_motor.rm_lin_flex[0] = 15.0 # m
rocket_motor.rm_lin_flex[1] =  5.0 # m
rocket_motor.rm_lin_flex[2] = 10.0 # m
rocket_motor.rm_lin_flex[3] = 20.0 # m

rocket_motor.rm_rot_flex[0] = trick.attach_units("degree", 15.0)
rocket_motor.rm_rot_flex[1] = trick.attach_units("degree",  3.0)
rocket_motor.rm_rot_flex[2] = trick.attach_units("degree",  9.0)
rocket_motor.rm_rot_flex[3] = trick.attach_units("degree",  6.0)
\end{code}

The unflexed thrust magnitude and the loss of thrust due to flexing is shown below:
\begin{longtable}{|r|r|r|}
\caption{Loss of thrust due to flexing}\\
\hline
\textbf{Sim time (s)} & \textbf{Thrust magnitude (N)} & \textbf{Delta thrust (N)}\\
\hline
\endhead

0.0 & 0.0 & 0.0 \\ \hline
0.1 & 2293.56 & -21.74 \\ \hline
0.2 & 15226.18 & -240.64 \\ \hline
0.3 & 49417.26 & -781.01 \\ \hline
0.4 & 44828.50 & -708.49 \\ \hline
0.5 & 37054.01 & -585.62 \\ \hline
0.6 & 23149.93 & -365.87 \\ \hline
0.7 & 5175.80 & -81.80 \\ \hline
0.8 & 1315.10 & -20.78 \\ \hline
0.9 & 450.89 & -7.13 \\ \hline
\end{longtable}

\subsection{Thrust Is Net}\label{sec:MultiNoz2}
When the thrust profile is a net thrust (\textbf{table\_is\_net\_thrust = True}), the thrust values represent the magnitude of the vector sum of the nozzle thrusts, not the sum of their vector magnitudes.  This scenario also assigns \textbf{cosine\_loss\_scale\_factor=0.9}.  Since 
$$Thrust_{net} = cosine\_loss\_scalefactor * Thrust_{raw}$$  the model must compute the pre-loss values in order to get an accurate mass consumption rate.  The thrust profile and pre-loss thrust values are below.
\begin{longtable}{|r|r|r|}
\caption{Net and pre-loss thrust values (straight cosine loss factor)}\\
\hline
\textbf{Sim time (s)} & \textbf{Net thrust (N)} & \textbf{Pre-loss thrust (N)}\\
\hline
\endhead

0.0 & 0.0 & 0.0 \\ \hline
0.1 & 2257.3 & 2508.1 \\ \hline
0.2 & 14986.5 & 16650.6 \\ \hline
0.3 & 48636.3 & 54040.3 \\ \hline
0.4 & 44120.0 & 49022.2 \\ \hline
0.5 & 36468.4 & 40520.4 \\ \hline
0.6 & 22784.1 & 25315.6 \\ \hline
0.7 & 5094.0 & 5660.0 \\ \hline
0.8 & 1294.3 & 1438.1 \\ \hline
0.9 & 450.9 & 501.0 \\ \hline
\end{longtable}

\subsection{Compute Cosine Losses}\label{sec:MultiNoz3}
This test will compute the cosine losses from the orientation and scale factor of the combined nozzles (as opposed to the straight loss calculation in \ref{sec:MultiNoz2}).  Note:  As in the previous run, \textbf{table\_is\_net\_thrust} must be set to True.

\begin{longtable}{|r|r|r|}
\caption{Net and pre-loss thrust values (computed cosine losses)}\\
\hline
\textbf{Sim time (s)} & \textbf{Net thrust (N)} & \textbf{Pre-loss thrust (N)}\\
\hline
\endhead

0.0 & 0.0 & 0.0 \\ \hline
0.1 & 2257.3 & 2540.0 \\ \hline
0.2 & 14986.5 & 16862.0 \\ \hline
0.3 & 48636.3 & 54726.3 \\ \hline
0.4 & 44120.0 & 49644.6 \\ \hline
0.5 & 36468.4 & 41034.9 \\ \hline
0.6 & 22784.1 & 25637.0 \\ \hline
0.7 & 5094.0 & 5731.9 \\ \hline
0.8 & 1294.3 & 1456.4 \\ \hline
0.9 & 450.9 & 507.3 \\ \hline
\end{longtable}

\subsection{Atmospheric Pressure}\label{sec:MultiNoz4}
This scenario tests the addition of the atmospheric pressure on the thrust calculations.  The run starts with an atmospheric pressure of zero, simulation a vacuum thrust.  At 0.3 seconds, a pressure of 10\textsuperscript{3} Pa is added.  As is shown below, the thrust drops.  At 0.5 seconds, an extreme pressure of 10\textsuperscript{10} Pa is introduced.  The thrust drops to zero.  A baseline run at vacuum is given for comparison.
\begin{figure}[H]
	\centering
	\includegraphics[width=\linewidth]{Images/AtmPressure01.png}
	\caption{Thrust variability due to atmospheric pressure}
\end{figure}


\subsection{Isp Table}\label{sec:MultiNoz5}
This unit test is specifically for seeing the changes in the data when the \textbf{mass\_flow\_rate} and Isp is calculated from the ISP propulsion consumption type profile data instead of the default \textbf{mburn} profile.
\begin{figure}[H]
	\centering
	\begin{subfigure}[b]{0.45\linewidth}
		\includegraphics[width=\linewidth]{Images/MassFlowRate01.png}
		\caption{Flow rate}
	\end{subfigure}
	\begin{subfigure}[b]{0.45\linewidth}
		\includegraphics[width=\linewidth]{Images/Isp01.png}
		\caption{Specific impulse}
	\end{subfigure}
	\caption{ISP vs mburn profile}
\end{figure}


\subsection{M-dot Table}\label{sec:MultiNoz6}
This unit test is specifically for seeing the changes in the data when the \textbf{mass\_flow\_rate} and Isp is calculated from the $\mathring{m}$ propulsion consumption type profile data instead of the default \textbf{mburn} profile.  The two plots overlap on each graph.
\begin{figure}[H]
	\centering
	\begin{subfigure}[b]{0.45\linewidth}
		\includegraphics[width=\linewidth]{Images/MassFlowRate02.png}
		\caption{Flow rate}
	\end{subfigure}
	\begin{subfigure}[b]{0.45\linewidth}
		\includegraphics[width=\linewidth]{Images/Isp02.png}
		\caption{Specific impulse}
	\end{subfigure}
	\caption{mdot vs mburn profile}
\end{figure}

\subsection{Small Tolerance}\label{sec:MultiNoz7}
This run tests the minimum operable magnitude of the vector \textbf{motor\_tolerance}.  If the magnitude of \textbf{motor\_tolerance} is less than \textbf{tolerance\_mag\_threshold}, consequential dispersion in attitude will be ignored and no calculations will be performed.  To illustrate, an unreasonably large dispersion threshold (1 rad) was chosen to force the test condition.  Below is the baseline z-component of thrust for the no threshold and threshold-applied scenarios.
\begin{figure}[ht!]
	\centering
	\includegraphics[width=\linewidth]{Images/SmallDispersion01.png}
	\caption{No threshold vs large dispersion error threshold}
\end{figure}

\subsection{Broken Nozzle}\label{sec:MultiNoz9}
In this scenario, at 0.5 seconds nozzle number four stops firing, and the remaining three nozzles increase their respective thrusts by a scaling factor of 1.1.  The input file looks like the following:
\begin{code}
trick.add_read(0.5,"""
rocket_motor.solid_motor.nozzle[0].sf_true*=1.1
rocket_motor.solid_motor.nozzle[1].sf_true*=1.1
rocket_motor.solid_motor.nozzle[2].sf_true*=1.1
rocket_motor.solid_motor.nozzle[3].shutdown_nozzle()
""")
\end{code}
With one engine shut down, the expectation is that the total moment will change.  Below illustrates  the effects with the x-component of moment.  The change in the x-component of moment is clearly seen at simulation time 0.5 seconds.
\begin{figure}[H]
	\centering
	\includegraphics[width=\linewidth]{Images/BrokenNozzle01.png}
	\caption{Nozzle shutdown}
\end{figure}


\subsection{Scale Factor Warnings}\label{sec:MultiNoz10}
This tests several scale factor error conditions.  A scale factor of zero generates this error:
\begin{figure}[ht!]
	\centering
	\includegraphics[width=\linewidth]{Images/ScaleFactorErrors01.png}
	\caption{Scale factor = 0}
\end{figure}

When $$cosine\_loss\_factor > 1$$ this error occurs:
\begin{figure}[H]
	\centering
	\includegraphics[width=\linewidth]{Images/ScaleFactorErrors02.png}
	\caption{Cosine loss greater than unity}
\end{figure}

When $$cosine\_loss\_factor < 0$$ this error occurs:
\begin{figure}[H]
	\centering
	\includegraphics[width=\linewidth]{Images/ScaleFactorErrors03.png}
	\caption{Negative cosine loss}
\end{figure}

\subsection{Bad Exit Area}\label{sec:MultiNoz11}
This unit simulation tests for a negative rocket nozzle area.
\begin{figure}[H]
	\centering
	\includegraphics[width=\linewidth]{Images/BadExitArea01.png}
	\caption{Negative nozzle area}
\end{figure}

\subsection{Consumption Type Error}\label{sec:MultiNoz12}
After one of the following has been chosen as the consumption type,
\begin{enumerate}
\item ConstantFlow
\item PropConsumptionIsp\item PropConsumptionMdot
\item PropConsumptionMburn
\end{enumerate}
it may not be switched to another type.
\begin{figure}[H]
	\centering
	\includegraphics[width=\linewidth]{Images/ConsumptionTypeChange01.png}
	\caption{Consumption type change}
\end{figure}

\subsection{Bad Update Call}\label{sec:MultiNoz13}
Update() is called prior to initialization, and update\_status() fails.  This does not cause the sim to abort.

\subsection{Duplicate Nozzle}\label{sec:MultiNoz14}
The model will not allow duplication of nozzles.  This runs tests that capability.
\begin{figure}[H]
	\centering
	\includegraphics[width=\linewidth]{Images/DuplicateNozzle01.png}
	\caption{Addition of duplicate nozzle}
\end{figure}

\subsection{Negative Thrust}\label{sec:MultiNoz15}
Negative thrust values are not allowed in thrust profile.
\begin{figure}[H]
	\centering
	\includegraphics[width=\linewidth]{Images/NegativeThrustValue01.png}
	\caption{Illegal thrust value}
\end{figure}

\subsection{Initialization Warnings (Fail)}\label{sec:MultiNoz20}
The rocket motor model does not allow re-initialization under normal conditions.  However, one may call the force\_initialize() method to allow this behavior.  This unit sim tests those conditions.
\begin{figure}[H]
	\centering
	\includegraphics[width=\linewidth]{Images/Reinitialize01.png}
	\caption{Model re-initialization}
\end{figure}

\subsection{NULL Flex, part A (Fail)}\label{sec:MultiNoz22a}
The linear flex perturbation must not be set to Null.
\begin{figure}[H]
	\centering
	\includegraphics[width=\linewidth]{Images/NullLinearFlex01.png}
	\caption{Null linear flex}
\end{figure}

\subsection{NULL Flex, part B (Fail)}\label{sec:MultiNoz22b}
The rotational flex perturbation must not be set to Null.
\begin{figure}[H]
	\centering
	\includegraphics[width=\linewidth]{Images/NullRotationalFlex01.png}
	\caption{Null rotational flex}
\end{figure}

\subsection{Flex Size Error (Fail)}\label{sec:MultiNoz23}
Flex array is too small.  The flex structure must be sized at 3 * number\_of\_nozzles.
\begin{figure}[H]
	\centering
	\includegraphics[width=\linewidth]{Images/FlexArray01.png}
	\caption{Flex array dimension too small}
\end{figure}

\subsection{No Nozzles (Fail)}\label{sec:MultiNoz24}
Must have at least one nozzle.  This insures some upward movement.
\begin{figure}[H]
	\centering
	\includegraphics[width=\linewidth]{Images/NoNozzle01.png}
	\caption{Must have at least one nozzle}
\end{figure}

\subsection{Nozzle Initialization, part A (Fail)}\label{sec:MultiNoz25a}
Transformation matrix must not have null pointer.
\begin{figure}[H]
	\centering
	\includegraphics[width=\linewidth]{Images/NullPointer01.png}
	\caption{No null pointer}
\end{figure}

\subsection{Nozzle Initialization, part B (Fail)}\label{sec:MultiNoz25b}
Position vector must not be null.
\begin{figure}[H]
	\centering
	\includegraphics[width=\linewidth]{Images/NoNullPosition01.png}
	\caption{No null position}
\end{figure}


\section{Thrust Table Verification}
The run scenarios that test the Thrust Table capability, \textbf{SIM\_table\_thrust\_rocket\_motor}, are as follows:
\begin{itemize}
  \item RUN\_mburn\_array (see \ref{sec:ThrustTable3})
  \item RUN\_mdot\_array (see \ref{sec:ThrustTable5})
  \item RUN\_isp\_array (see \ref{sec:ThrustTable1})    
  \item RUN\_mburn\_vector (see \ref{sec:ThrustTable4})
  \item RUN\_mdot\_vector (see \ref{sec:ThrustTable6})
  \item RUN\_same\_consumption\_type (see \ref{sec:ThrustTable7})    
  \item RUN\_isp\_vector (see \ref{sec:ThrustTable2})
\end{itemize}


\subsection{Mburn Array}\label{sec:ThrustTable3}
This unit test is used as a baseline.  It also checks the disable\_flex() functionality.  The data for these unit simulations can be found in the class \textbf{RocketMotorArrayData}, in the \textbf{example\_rocket\_motor.hh} file.  Below is the thrust profile and mass consumption profile used in this and the following scenarios:
\begin{figure}[H]
	\centering
	\begin{subfigure}[b]{0.45\linewidth}
		\includegraphics[width=\linewidth]{Images/MburnArray01.png}
		\caption{Thrust profile}
	\end{subfigure}
	\begin{subfigure}[b]{0.45\linewidth}
		\includegraphics[width=\linewidth]{Images/MburnArray02.png}
		\caption{Consumption profile}
	\end{subfigure}
	\caption{Thrust and Mass consumption}
\end{figure}

\subsection{Mdot Array}\label{sec:ThrustTable5}
This scenario uses the \textbf{RUN\_mburn\_array} configuration.  In order to configure the $\mathring{m}$ run, the model computes the derivative of the mburn array to get the $\mathring{m}$ array values.

\subsection{Isp Array}\label{sec:ThrustTable1}
This scenario generates Isp values from a pre-defined time series array.  A specific example can be found in \textbf{example\_rocket\_motor.hh}.

\subsection{Mburn Vector}\label{sec:ThrustTable4}
Identical to \ref{sec:ThrustTable3}, but uses C++ STL vectors.

\subsection{Mdot Vector}\label{sec:ThrustTable6}
Uses \ref{sec:ThrustTable4} as the data source (vectors rather than arrays) for constructing the $\mathring{m}$ array values.

\subsection{Same Consumption Type}\label{sec:ThrustTable7}
At 0.3 simulation time, this run switches to the same consumption type:
\begin{code}
set_consumption_type(
        trick.RocketMotor_TableThrust.PropConsumptionMburn)
\end{code}

There is no change in the simulation results.

\subsection{Isp Vector}\label{sec:ThrustTable2}
Identical to \ref{sec:ThrustTable1}, but uses C++ STL vectors.

\end{document}